% =============================================================================
%  Lab 0 - VHDL Fundamentals   *** MODEL CHAPTER ***
%
%  This chapter is written out in full as an example of the standard expected.
%  Read it once. Then back-fill your own Lab 0 work in the same shape.
%
%  Notice what it does: every table has numbered rows, every equation names its
%  simplification law, every figure has a sentence saying what it proves, and
%  the Discussion describes a real mistake.
% =============================================================================

\chapter{Lab 0 --- VHDL Fundamentals}

\section{Objective}

Introduce the structure of a VHDL design --- the \texttt{entity} that declares an
interface and the \texttt{architecture} that implements it --- and verify a small
combinational circuit by simulation rather than by inspection. The circuit used
here is a majority-of-three detector.

\section{Specification}

The output $Y$ is asserted when at least two of the three inputs are asserted.
With three binary inputs the table is small enough to enumerate completely, so it
is given exhaustively (Type E).

\begin{table}[htbp]
\centering
\caption{Exhaustive truth table (Type E) for \texttt{majority3}.}
\label{tab:lab00-truth}
\begin{tabular}{cccc||cl}
\toprule
\# & $A$ & $B$ & $C$ & $Y$ & minterm \\
\midrule
0 & 0 & 0 & 0 & 0 & \\
1 & 0 & 0 & 1 & 0 & \\
2 & 0 & 1 & 0 & 0 & \\
3 & 0 & 1 & 1 & 1 & $m_3$ \\
4 & 1 & 0 & 0 & 0 & \\
5 & 1 & 0 & 1 & 1 & $m_5$ \\
6 & 1 & 1 & 0 & 1 & $m_6$ \\
7 & 1 & 1 & 1 & 1 & $m_7$ \\
\bottomrule
\end{tabular}
\end{table}

\subsection*{Boolean derivation}

\begin{align*}
\text{(1) canonical} \quad & Y = \textstyle\sum m(3,5,6,7)
   = \overline{A}BC + A\overline{B}C + AB\overline{C} + ABC \\
\text{(2) simplified} \quad & Y = AB + AC + BC
   \qquad \text{[$ABC$ is adjacent to each of the other three minterms, so each}\\
   & \qquad\qquad\qquad\qquad\quad\ \text{pair combines and the $C$, $B$, $A$ literals drop out]}\\
\text{(3) VHDL} \quad & \texttt{Y <= (A and B) or (A and C) or (B and C);}
\end{align*}

Note that $Y$ does not appear on the right-hand side of its own equation. This
circuit has no memory: its output depends only on the inputs present at that
instant.

\section{Design}

The entity declares four ports: three \texttt{std\_logic} inputs and one
\texttt{std\_logic} output. A behavioural architecture is appropriate here
because the simplified equation maps directly onto a single concurrent signal
assignment; describing it structurally as three AND gates and an OR gate would
add three component declarations without changing the synthesised result.

\section{Implementation}

\vhdlfile{code/lab00/majority3.vhd}{\texttt{majority3} --- majority-of-three detector}

\section{Verification}

All eight input combinations were applied. When a truth table is small enough to
enumerate exhaustively, exhaustive testing is the coverage argument: no input
combination exists that was not checked, so no untested behaviour can remain.

The testbench is driven from an array of vectors rather than from hand-written
blocks, so that the test data and the truth table are visibly the same object.

\section{Results}

\figplaceholder{the RTL Viewer output for \texttt{majority3}}
% Example of the sentence required with each figure:
% "The RTL Viewer shows two levels of combinational logic and no storage
%  elements, confirming the design is purely combinational as specified."

\figplaceholder{the Questa waveform showing all eight input combinations}
% "Y is high only during rows 3, 5, 6 and 7, matching Table 1."

Compilation completed with 0 errors and 0 warnings. The simulation transcript
reported \texttt{PASS --- majority3: all 8 rows correct}, so every assertion in
the testbench held.

\section{Discussion}

The first version of the testbench applied all eight vectors but checked none of
them: the assignments were present and the \texttt{assert} statements were not.
The waveform looked correct, so the design appeared to be verified when in fact
nothing had been verified at all.

The lesson generalises, and is the reason every testbench in this report is
self-checking: \emph{a waveform is evidence, an assertion is proof}. Looking at a
waveform and deciding it seems right is not verification, because it only tests
whether the output matches what the person looking at it already expected.
